\documentclass[../master.tex]{subfiles}

\begin{document}
\chapter{スタビライザー形式}
これまではShor符号の具体的な構成を通して、
量子誤り訂正符号の基本的なアイデアを紹介してきた。
このアプローチは、
情報を保護する量子状態を具体的にどのような形で符号化するかを先に定め、
その符号化を実現する回路を構成し、
その後に符号化された状態を壊さずに誤りを検出するための物理量を見つけるというものであった。
しかし、より高い誤り耐性などの性能を求めて複雑な符号を構築しようとする際、
この方法は非常に煩雑になり、いずれ限界を迎えることになる。
古典線形符号において符号は生成行列から作るものという視点だけではなく、
検査行列の核$(\ker H)$のなかの元という視点でも特徴づけることができる。
これと同様に量子誤り訂正においても、まずシンドローム測定のための物理量を先に与え、
その満たすべき条件から論理状態のなす空間を定めるという立場も取ることができる。
本章で扱うスタビライザー形式はこの考えを一般化したものであり、
可換な一般化パウリ演算子の集合を用いて符号空間をその共通固有空間として記述する
\cite{Gottesman1997-ep,Daniel-Gottesman1999-eh,Bullock2007-jl}。
この枠組みによって、
エラーの検出及び訂正に関わるシンドローム測定、
論理状態間を区別する論理演算子の構造が整理され、
さらに符号空間を可換ハミルトニアンの基底状態のなす空間としてみる見方も得られる。

\section{符号語から検査条件へ}
これまでの Shor 符号の議論では、まず論理状態の具体的な形を与え、
その符号化を実現する符号化する回路を構成し、
その後に誤りを検出するためのシンドローム測定を考えてきた。
このような構成は量子誤り訂正の基本的な仕組みを理解する上では有用である。
実際、どのような量子状態が保護されるべきか、
またどのような測定がその状態を壊さずに誤りの情報だけを読み出すのかを、
具体例を通して直接確かめることができるからである。

しかし、符号距離などの性能や実装との相性を考慮するためより一般の量子誤り訂正符号を構成しようとすると、
このように符号語や回路を先に与える立場は急速に煩雑になる。
例えばCSS符号では、シフトエラーと位相エラーの両方の検出を同時に行える条件は、
符号の検査行列の関係式であるCSS条件 $H_sH_p^\mathsf{T}=0$ を満たすことである。
これでは符号語から始めた場合、生成行列$G$に対応する符号化を考えて、
そこから誤りを検出するための$H_s,H_p$を読み取り、
それがCSS条件を満たすかどうかを確認するというようになる。
このように符号語から始めるアプローチの場合、
量子誤り訂正符号としてそもそも使えるかどうかを確認するのでさえ、非常に遠回りになっており、
そこから性能を考慮して符号を改良することもさらに難しくなってしまう。

ここで古典線形符号の場合を思い出そう。
古典線形符号では、符号語のなす空間(符号空間)は生成行列 $G$ の像$\Im G$ として与えることができる一方で、
検査行列 $H$ の核 $\ker H$ としても与えることができる。
すなわち、符号はメッセージからどのように符号語を作るかという立場から見ることも、
どのような検査条件を満たすものを符号語とみなすかという立場から見ることもできる。
後者の立場では、まず $H$ を与えることで符号語空間そのものが定まり、
その空間の基底を選ぶことによって対応する生成行列 $G$ を構成できる。
したがって、生成行列による記述と検査行列による記述は、
同じ符号空間の2つの異なる見方にほかならない。

量子誤り訂正符号においても同様の発想をとることができる。
すなわち、先に論理状態や符号化の回路を与えるのではなく、
まず誤りのシンドロームを与える可換な検査演算子を定め、
それらを満たす状態全体を符号空間とみなすのである。
この立場をとれば、シンドローム測定の回路は
既に与えられた量子状態を調べるための補助的な装置ではなく、
むしろ符号空間そのものを定義するものとして理解される。

\section{スタビライザー群と符号空間}
ではシンドローム測定で測定する物理量をどのように与えるかを考えよう。
$\mathbb{F}_p$におけるHadamardテストでは、
ユニタリ演算子の固有値が$\omega^k = \exp(\frac{2\pi i}{p}k)$であれば、
測定するquditは確実に状態$\ket{-k}$になることを見た。
このことは、誤りのシンドロームを測定するための物理量は、
誤りのシンドロームを表すユニタリ演算子の固有値が$\omega^k$であれば、
誤りの検出には都合が良いことを示唆している。
このような固有値を取るようなユニタリ演算子は一般化パウリ演算子の組み合わせで構成される。
そのような演算子の集合を考える必要があるのとわかる。

さらに、複数のシンドロームを同時に測定したいのであれば、
対応する検査演算子は測定する順番によって結果が変わらないようにしなければならない。
すなわち、2つの検査演算子 $S_1,S_2$ に対して、
\begin{equation}\begin{aligned}[b]
    S_1S_2\ket{\psi} &= S_2S_1\ket{\psi}  &
    \Rightarrow \quad S_1S_2S_1^{-1}S_2^{-1}&=I
\end{aligned}\end{equation}
が成り立っていてほしい。
この式には演算子の積、単位元、逆元が現れており、
シンドローム測定に用いる演算子の集合を整理するには、
これらを備えた代数的枠組みとして群を用いるのがよいであろう。
そうした一般化パウリ演算子とその積のなす群が一般化パウリ群である。

\begin{NoteBox}[title={一般化パウリ群}]{}
    $\mathbf a=(a_1,\dots,a_n),\, \mathbf b=(b_1,\dots,b_n)\in\mathbb F_p^n$
    に対して
    \begin{equation}
        \mathcal X^{\mathbf a}:=\mathcal X_1^{a_1}\cdots \mathcal X_n^{a_n},\qquad
        \mathcal Z^{\mathbf b}:=\mathcal Z_1^{b_1}\cdots \mathcal Z_n^{b_n}
    \end{equation}
    と定める。
    一般化パウリ群$\mathcal{P}_p^n$とは、$n$ 個の qudit に作用する一般化パウリ演算子全体のなす群で、
    この群の元は
    \begin{equation}\begin{aligned}[b]
        \mathcal{P}_p^n=
        \qty{
            \omega^k \mathcal{X}^{\mathbf{a}}\mathcal{Z}^{\mathbf{b}}
            \mid k\in  Z_p,  \mathbf{a},\mathbf{b}\in \mathbb{F}_p^n }
    \end{aligned}\end{equation}
    と書け、積は
    \begin{equation}\begin{aligned}[b]
        (\omega^{k_1} \mathcal{X}^{\mathbf{a}_1}\mathcal{Z}^{\mathbf{b}_1})\cdot
        (\omega^{k_2} \mathcal{X}^{\mathbf{a}_2}\mathcal{Z}^{\mathbf{b}_2})
        &=
        \omega^{k_1+k_2+\mathbf{b}_1\cdot\mathbf{a}_2}
        \mathcal{X}^{\mathbf{a}_1+\mathbf{a}_2}\mathcal{Z}^{\mathbf{b}_1+\mathbf{b}_2}
    \end{aligned}\end{equation}
    のように定まる。
\end{NoteBox}

このうち、シンドローム測定に用いる検査演算子は、
測定する順番によって結果が変わらないように、互いに可換でなければならない。
したがって、一般化パウリ群の中でも可換な部分群が検査演算子の集合となる。

検査演算子の性質として他にも重要な要請がある。
符号語となる状態$\ket{\psi}_L$に対しては各検査演算子がいつも決まった固有値を与えるようにしたい。
すなわち、ある検査演算子 $S$ に対して
\begin{equation}
   S\ket{\psi}_L=\omega^k\ket{\psi}_L
\end{equation}
となるとする。
ここで固有値 $\omega^k$ 自体は一般には $1$ である必要はない。
しかし、$\omega^{-k}S$ もまた一般化パウリ演算子であるから、
検査演算子を位相倍だけ取り直せば
\begin{equation}
    (\omega^{-k}S)\ket{\psi}_L=(+1)\ket{\psi}_L
\end{equation}
と書くことができる。
したがって、各検査演算子の固有値は符号状態に対して $+1$ であるとする。

また、この要請のもとでは検査演算子の集合の中に非自明な位相倍の単位演算子 $\omega^k I$ ($k\neq 0$) を含めることはできない。
なぜなら、
\begin{equation}
    \omega^k I\ket{\psi}_L=\ket{\psi}_L
\end{equation}
を満たす状態 $\ket{\psi}_L$ は$\ket{\psi}_L=0$であり計算に使う状態としては不適当だからである。
したがって、符号空間を非自明に保つためには、$\omega^kI (k\neq 0)$を除かなければならない。

以上の要請をまとめると、シンドローム測定に用いる検査演算子としては、
一般化パウリ群の中で互いに可換であり、かつ非自明な位相倍の単位演算子を含まない部分群である。
このような部分群をスタビライザー群と呼び、
いままでシンドローム測定で測定されてきた検査演算子のことをスタビライザー演算子と呼ぶ。
\begin{NoteBox}[title={スタビライザー群}]{}
スタビライザー群とは、次の2つの条件を満たす一般化パウリ群 $\mathcal{P}_{p}^n$ の部分群 $\mathcal{S}$ のことである。
    \begin{itemize}
        \item $\mathcal{S}\subset \mathcal{P}_{p}^n$ は可換な部分群である。
        \item 任意の $k\neq 0$ に対して $\omega^k I \notin \mathcal{S}$.
    \end{itemize}
\end{NoteBox}
このスタビライザー群の定義から、具体的にどのような演算子がスタビライザー群の元になりうるかを考えてみよう。
スタビライザー群の元は一般化パウリ群の元であるから、ある $k\in Z_p$ と $\mathbf{a},\mathbf{b}\in\mathbb{F}_p^n$ に対して
\begin{equation}
    S=\omega^k \mathcal{X}^{\mathbf{a}}\mathcal{Z}^{\mathbf{b}}
\end{equation}
の形で書ける。
スタビライザー群の定義から、スタビライザー群の元は互いに可換でなければならない。
すなわち、スタビライザー群の元 $S_1=\omega^{k_1} \mathcal{X}^{\mathbf{a}_1}\mathcal{Z}^{\mathbf{b}_1}$ と
$S_2=\omega^{k_2} \mathcal{X}^{\mathbf{a}_2}\mathcal{Z}^{\mathbf{b}_2}$ に対して以下が成り立つ。
\begin{equation}\begin{aligned}[b]
    0&=\qty(\omega^{k_1} \mathcal{X}^{\mathbf{a}_1}\mathcal{Z}^{\mathbf{b}_1})\cdot
    \qty(\omega^{k_2} \mathcal{X}^{\mathbf{a}_2}\mathcal{Z}^{\mathbf{b}_2})
    -\qty(\omega^{k_2} \mathcal{X}^{\mathbf{a}_2}\mathcal{Z}^{\mathbf{b}_2})\cdot
    \qty(\omega^{k_1} \mathcal{X}^{\mathbf{a}_1}\mathcal{Z}^{\mathbf{b}_1})\\
    &=
    \omega^{k_1+k_2+\mathbf{b}_1\cdot\mathbf{a}_2}
    \mathcal{X}^{\mathbf{a}_1+\mathbf{a}_2}\mathcal{Z}^{\mathbf{b}_1+\mathbf{b}_2}
    -
    \omega^{k_1+k_2+\mathbf{b}_2\cdot\mathbf{a}_1}
    \mathcal{X}^{\mathbf{a}_1+\mathbf{a}_2}\mathcal{Z}^{\mathbf{b}_1+\mathbf{b}_2}\\
    &=\qty(\omega^{\mathbf{b}_1\cdot\mathbf{a}_2-\mathbf{b}_2\cdot\mathbf{a}_1}-1)
    \omega^{k_1+k_2+\mathbf{b}_2\cdot\mathbf{a}_1}
    \mathcal{X}^{\mathbf{a}_1+\mathbf{a}_2}\mathcal{Z}^{\mathbf{b}_1+\mathbf{b}_2}.
\end{aligned}\end{equation}
この式が成り立つためには、
\begin{equation}
    \mathbf{b}_1\cdot\mathbf{a}_2-\mathbf{b}_2\cdot\mathbf{a}_1=0
\end{equation}
でなければならない。
これをシンプレクティック条件と言い、スタビライザー群の元同士が可換である条件である。
さらに、スタビライザー群の定義から、$\omega^k I$ ($k\neq 0$) をスタビライザー群の元として含めることはできない。
したがって、スタビライザー群を生成する一般化パウリ演算子
\begin{equation}
    S_i=\omega^{k_i}\mathcal X^{\mathbf a_i}\mathcal Z^{\mathbf b_i}
\end{equation}
のラベル $(\mathbf a_i,\mathbf b_i)$ は、
互いにこのシンプレクティック条件を満たさなければならない。

CSS 符号もスタビライザー符号の一種であることを確認してみよう。
CSS 符号の検査行列は $H_s$ と $H_p$ であるから、
CSS 符号のスタビライザー演算子は、シフトエラーに対応するものと位相エラーに対応するものへと分かれる。
それぞれ各検査行列の行を $\mathbf{b}_i=H_s[i,:]$ および $\mathbf{a}_i=H_p[i,:]$ とおくと、
スタビライザー演算子は
\begin{equation}\begin{aligned}[b]
    S_{si} &= \mathcal{Z}^{\mathbf b_i} = \prod_m \mathcal{Z}_m^{H_s[i,m]} &
    S_{pi} &= \mathcal{X}^{\mathbf a_i} = \prod_m \mathcal{X}_m^{H_p[i,m]}
\end{aligned}\end{equation}
の形で書ける。
このとき、シンプレクティック条件は $H_sH_p^\mathsf{T}=0$ と同値である。
というのも、CSS 型ではシフトエラーに対応する演算子に対しては $\mathbf{a}_{si}=0$,
位相エラーに対応する演算子に対しては $\mathbf{b}_{pj}=0$ であるから、
シンプレクティック条件は
\begin{equation}\begin{aligned}[b]
    0&=\mathbf{b}_{si}\cdot\mathbf{a}_{pj}-\mathbf{b}_{pj}\cdot\mathbf{a}_{si}
    =-H_s[i,:]\cdot H_p[j,:]^\mathsf{T}
\end{aligned}\end{equation}
と書けるからである。
また、このようなCSS型のスタビライザーでは$\mathbf{a}_{si}=\mathbf{b}_{pj}=0$であるから、
スタビライザーは $S_{pi}=\omega^{k_i}\mathcal{X}^{\mathbf a_i}$ や $S_{si}=\omega^{k_i}\mathcal{Z}^{\mathbf b_i}$ の形で書ける。
さらに前に述べたように、スタビライザー演算子は位相倍だけ取り直して符号空間上で固有値が $+1$ になるように選んでよい。
したがって、CSS型のスタビライザー演算子は代表元の取り方として
$S_{pi}=\mathcal{X}^{\mathbf a_i}$ や $S_{si}=\mathcal{Z}^{\mathbf b_i}$ の形にしてよい。

また、スタビライザー群の各元が固有値 $+1$ を与える状態全体を符号空間と考えることができる。
この符号空間のことをスタビライザー符号と呼ぶ。
\begin{NoteBox}[title={スタビライザー符号}]{}
    $\mathcal{S}$ の全ての元に対して固有値 $+1$ をもつ状態全体
    \begin{equation}
        \mathcal{C}
        =\qty{
            \ket{\psi}\in (\mathbb{C}^p)^{\otimes n}
            \mid S\ket{\psi}=\ket{\psi},\ \forall S\in\mathcal{S}}
    \end{equation}
    をスタビライザー符号の符号空間という。
\end{NoteBox}

この定義をShor符号で確認してみよう。
前章で見た qudit 版 Shor符号の論理状態は
\begin{equation}
    \ket{j}_L = \ket{\tilde{j}}_l \ket{\tilde{j}}_l \ket{\tilde{j}}_l
\end{equation}
であり、各ブロックの内部ではシフトエラーを検出するために
\begin{equation}
    \mathcal{Z}_1\mathcal{Z}_2^\dagger,\ \mathcal{Z}_2\mathcal{Z}_3^\dagger,\
    \mathcal{Z}_4\mathcal{Z}_5^\dagger,\ \mathcal{Z}_5\mathcal{Z}_6^\dagger,\
    \mathcal{Z}_7\mathcal{Z}_8^\dagger,\ \mathcal{Z}_8\mathcal{Z}_9^\dagger
\end{equation}
を測定し、ブロック間では位相エラーを検出するために
\begin{equation}
    \mathcal{X}_1\mathcal{X}_2\mathcal{X}_3
    \mathcal{X}_4^\dagger\mathcal{X}_5^\dagger\mathcal{X}_6^\dagger,\qquad
    \mathcal{X}_4\mathcal{X}_5\mathcal{X}_6
    \mathcal{X}_7^\dagger\mathcal{X}_8^\dagger\mathcal{X}_9^\dagger
\end{equation}
を測定していた。
これらはいずれも論理状態に対して固有値 $+1$ を与える。
例えば、
\begin{equation}\begin{aligned}[b]
    (\mathcal{Z}_1\mathcal{Z}_2^\dagger)\ket{j}_L &= \ket{j}_L, &
    (\mathcal{Z}_2\mathcal{Z}_3^\dagger)\ket{j}_L &= \ket{j}_L, \\
    (\mathcal{X}_1\mathcal{X}_2\mathcal{X}_3
    \mathcal{X}_4^\dagger\mathcal{X}_5^\dagger\mathcal{X}_6^\dagger)\ket{j}_L &= \ket{j}_L
\end{aligned}\end{equation}
となる。
したがって、Shor符号の符号空間はこれらの演算子すべての固有値 $+1$ の共通固有空間として与えられる。

ここで重要なのは、スタビライザー群とはこれら8個の演算子そのものではなく、
それらの積も含めた群全体であるという点である。
例えば
\begin{equation}
    (\mathcal{Z}_1\mathcal{Z}_2^\dagger)(\mathcal{Z}_2\mathcal{Z}_3^\dagger)
    =\mathcal{Z}_1\mathcal{Z}_3^\dagger
\end{equation}
もスタビライザー群の元であり、やはり符号空間上では固有値 $+1$ を与える。
これは古典線形符号において、符号空間 $\ker H$ を定めるのに検査行列の全ての行が必要なわけではなく、
独立な行だけを取れば十分であったことに対応している。
例えば、3ビット繰り返し符号 $0\to 000,\ 1\to 111$ では、
3つのビットを全て比較する検査条件として
\begin{equation}
    H=
    \begin{pmatrix}
        1 & -1 & 0\\
        0 & 1 & -1\\
        1 & 0 & -1
    \end{pmatrix}
\end{equation}
を考えることもできるが、3行目は前2行の和で表されるため冗長であり、
$\ker H$ を定めるには独立な2行だけで十分である。
スタビライザー群でも同様に、符号空間を定めるには独立な検査条件だけを与えれば十分である。
一方で、スタビライザー群そのものは、それらの検査演算子の積で得られる元も含めた可換部分群として定義される。
したがって、積で得られる元は新しい独立な検査条件を加えているのではなく、
もとの検査条件から従うものを群として閉じた形でまとめたものとみなせる。

このことは、スタビライザー群を記述するのに群の全ての元を逐一並べる必要はなく、
いくつかの基本的な演算子を指定すれば十分であることを意味している。
実際、それらの演算子から積を作っていくことで、スタビライザー群の全ての元を得ることができる。
このように群全体を作るもとになる演算子を生成元と呼ぶ。

独立な生成元を $S_1,S_2,\dots,S_r$ とすると、
それらから作られるスタビライザー群は
\begin{equation}
    \mathcal{S}=\langle S_1,S_2,\dots,S_r\rangle
\end{equation}
と書かれる。
ここで $\langle S_1,S_2,\dots,S_r\rangle$ は、
$S_1,S_2,\dots,S_r$ の積によって得られる全ての元からなる群を表している。
このとき、符号空間は
\begin{equation}
    S_i\ket{\psi}=\ket{\psi}\qquad (i=1,\dots,r)
\end{equation}
を全て満たす状態全体として記述できる。
したがって、実際のシンドローム測定においても、
群の全ての元を測定する代わりに、生成元に対応する検査演算子だけを測定すれば十分である。


\section{論理演算子}
ここまでで、スタビライザー群を与えることによって符号空間そのものが定まることを見た。
しかし、量子誤り訂正符号として情報を保持するためには、
どの状態が符号空間に属するかが分かるだけでは十分でない。
前章で扱った qudit 版 Shor符号では、符号化された状態を $\ket{j}_L$ と書いていた。
ところが、スタビライザー群を与えただけでは、
符号空間の中でこのラベル $j$ をどのように区別するのかはまだ見えていない。
そこで次には、符号空間を保ちながら、
論理状態 $\ket{j}_L$ に応じた位相を与えたり、
別の論理状態 $\ket{j'}_L$ へ移したりする演算子を考えたい。
このような演算子を論理演算子と呼ぶ。

このことを qudit 版 Shor符号で具体的に見てみよう。
前章で導入した論理状態は
\begin{equation}\begin{aligned}[b]
    \ket{j}_L &= \ket{\tilde{j}}_l \ket{\tilde{j}}_l \ket{\tilde{j}}_l &
    \ket{\tilde{j}}_l &= \frac{1}{\sqrt{p}}\sum_{k=0}^{p-1}\omega^{-jk}\ket{kkk}
\end{aligned}\end{equation}
であった。
物理的な一般化パウリ演算子 $\mathcal{X}$ は状態のラベルを1つシフトする演算子であり、
$\mathcal{Z}$ は状態のラベルに応じた位相を与える演算子である。
したがって、論理演算子についても
\begin{equation}
    \overline{\mathcal{X}}\ket{j}_L=\ket{j+1}_L,\qquad
    \overline{\mathcal{Z}}\ket{j}_L=\omega^j\ket{j}_L
\end{equation}
を満たしていてほしい。
ただし、ここで用いている論理状態は $\ket{\tilde{j}}$ 基底で書かれているため、
物理的な $\mathcal{X}$ と $\mathcal{Z}$ の役割が見かけ上入れ替わることに注意が必要である。
この論理状態のラベル $j$ を読み出す演算子として、
\begin{equation}
    \overline{\mathcal{Z}} := \mathcal{X}_1\mathcal{X}_2\mathcal{X}_3
\end{equation}
を考える。
実際、
\begin{equation}\begin{aligned}[b]
    \overline{\mathcal{Z}}\ket{j}_L
    = (\mathcal{X}_1\mathcal{X}_2\mathcal{X}_3)
    \qty(\ket{\tilde{j}}_l \ket{\tilde{j}}_l \ket{\tilde{j}}_l)
    = \qty(\mathcal{X}_1\mathcal{X}_2\mathcal{X}_3\ket{\tilde{j}}_l)
    \ket{\tilde{j}}_l \ket{\tilde{j}}_l
    = \omega^{j}\ket{j}_L
\end{aligned}\end{equation}
となる。
したがって、$\overline{\mathcal{Z}}$ は論理状態 $\ket{j}_L$ のラベル $j$ に応じた位相を与える演算子として働く。

一方、論理状態のラベルを1つ進める演算子として
\begin{equation}
    \overline{\mathcal{X}} := \mathcal{Z}_1^\dagger \mathcal{Z}_4^\dagger \mathcal{Z}_7^\dagger
\end{equation}
を考える。
これは、各ブロックに対して同じだけラベルをずらしながら
$\ket{\tilde{j}}_l\ket{\tilde{j}}_l\ket{\tilde{j}}_l$ という形を保つような最も単純な候補である。
$\mathcal{Z}^\dagger$ は $\ket{\tilde{j}}_l$に作用させると
\begin{equation}\begin{aligned}[b]
    \mathcal{Z}_1^\dagger\ket{\tilde{j}}_l
    = \frac{1}{\sqrt{p}}\sum_{k=0}^{p-1}\omega^{-jk}\mathcal{Z}_1^\dagger\ket{kkk}_{123}
    = \frac{1}{\sqrt{p}}\sum_{k=0}^{p-1}\omega^{-(j+1)k}\ket{kkk}_{123}
    = \ket{\widetilde{j+1}}_l
\end{aligned}\end{equation}
となるので、
\begin{equation}\begin{aligned}[b]
    \overline{\mathcal{X}}\ket{j}_L
    = (\mathcal{Z}_1^\dagger \mathcal{Z}_4^\dagger \mathcal{Z}_7^\dagger)
    \qty(\ket{\tilde{j}}_l \ket{\tilde{j}}_l \ket{\tilde{j}}_l)
    = \ket{\widetilde{j+1}}_l \ket{\widetilde{j+1}}_l \ket{\widetilde{j+1}}_l
    = \ket{j+1}_L
\end{aligned}\end{equation}
となる。
したがって、$\overline{\mathcal{X}}$ は論理状態のラベルを1つずらす演算子となっている。

また、これらの演算子はスタビライザー演算子と可換である。
例えば
\begin{equation}\begin{aligned}[b]
    (\mathcal{Z}_1\mathcal{Z}_2^\dagger)(\mathcal{X}_1\mathcal{X}_2\mathcal{X}_3)
    = (\mathcal{Z}_1\mathcal{X}_1)(\mathcal{Z}_2^\dagger\mathcal{X}_2)\mathcal{X}_3
    = (\omega\mathcal{X}_1\mathcal{Z}_1)(\omega^{-1}\mathcal{X}_2\mathcal{Z}_2^\dagger)\mathcal{X}_3
    = (\mathcal{X}_1\mathcal{X}_2\mathcal{X}_3)(\mathcal{Z}_1\mathcal{Z}_2^\dagger)
\end{aligned}\end{equation}
であり、位相因子が打ち消し合うとわかる。
他のスタビライザーに対しても同様に可換性が成り立つ。
一方で、$\overline{\mathcal{X}}$ や $\overline{\mathcal{Z}}$ はスタビライザー群の元ではない。
実際、$\overline{\mathcal{Z}}=\mathcal{X}_1\mathcal{X}_2\mathcal{X}_3$ は1つのブロックのみに作用する3体演算子であるが、
$\mathcal{X}$ 型のスタビライザー生成元は2つの隣接するブロックにまたがる6体演算子であるため、
それらの積からこの形を作ることはできない。
また、$\overline{\mathcal{X}}=\mathcal{Z}_1^\dagger \mathcal{Z}_4^\dagger \mathcal{Z}_7^\dagger$ は
各ブロックから1サイトずつを選んだ形で作用するが、
$\mathcal{Z}$ 型のスタビライザー生成元の積から得られるのは各ブロック内部での差分を表す演算子のみであり、
やはりこの形は得られない。
さらに、これら2つの演算子は
\begin{equation}\begin{aligned}[b]
    \overline{\mathcal{Z}}\overline{\mathcal{X}}
    &= (\mathcal{X}_1\mathcal{X}_2\mathcal{X}_3)
    (\mathcal{Z}_1^\dagger \mathcal{Z}_4^\dagger \mathcal{Z}_7^\dagger)\\
    &= \omega
    (\mathcal{Z}_1^\dagger \mathcal{Z}_4^\dagger \mathcal{Z}_7^\dagger)
    (\mathcal{X}_1\mathcal{X}_2\mathcal{X}_3)
    = \omega\,\overline{\mathcal{X}}\overline{\mathcal{Z}}
\end{aligned}\end{equation}
を満たす。
したがって、$\overline{\mathcal{X}}$ と $\overline{\mathcal{Z}}$ は論理 qudit 上の一般化パウリ演算子として期待される交換関係を持っている。
これらは符号空間を保ちながらその内部で非自明に作用する論理演算子の具体例となっている。

ここで、一般に論理演算子が満たすべき条件を整理しておこう。
一般に、論理演算子の候補を一般化パウリ群の元 $P$ として考える。
論理演算子は符号空間を保つ必要があるので、任意の $\ket{\psi}\in\mathcal{C}$ に対して
\begin{equation}
    P\ket{\psi}\in\mathcal{C}
\end{equation}
を満たさなければならない。
ここで $\ket{\psi}\in\mathcal{C}$ は全てのスタビライザー演算子 $S\in\mathcal{S}$ に対して
\begin{equation}
    S\ket{\psi}=\ket{\psi}
\end{equation}
を満たしていたから、$P\ket{\psi}$ も符号空間に属するためには
\begin{equation}
    S(P\ket{\psi})=P\ket{\psi}\qquad (\forall S\in\mathcal{S})
\end{equation}
が成り立たなければならない。
これは
\begin{equation}
    P^{-1}SP\ket{\psi}=\ket{\psi}\qquad (\forall S\in\mathcal{S})
\end{equation}
と書き換えられる。
したがって、$P$ が符号空間を保つためには、任意の $S\in\mathcal{S}$ に対して
$P^{-1}SP$ も再び $\mathcal{S}$ に属していなければならない。
一方で、スタビライザー群の元そのものは符号空間上で恒等的に作用するため、
論理状態を区別することはできない。
したがって、求めるべき論理演算子とは、
符号空間を保ちながらも、
スタビライザー群の元としては自明でないような演算子である。
このような論理演算子全体を系統的に特徴づけるために、正規化群を導入する。

この条件を満たす一般化パウリ演算子全体は群をなし、
スタビライザー群 $\mathcal{S}$ の正規化群と呼ばれる。
\begin{NoteBox}[title={正規化群}]{}
    スタビライザー群 $\mathcal{S}\subset\mathcal{P}_p^n$ の正規化群とは、
    \begin{equation}
        N(\mathcal{S})
        :=
        \qty{
            P\in\mathcal{P}_p^n
            \mid
            P^{-1}SP\in\mathcal{S},\ \forall S\in\mathcal{S}}
    \end{equation}
    で定まる部分群である。
\end{NoteBox}
この集合が実際に群となることは、一般化パウリ群の部分集合として比較的容易に確かめられる。
実際、単位演算子 $I$ は自明に $N(\mathcal{S})$ に属する。
また、$P,Q\in N(\mathcal{S})$ ならば、任意の $S\in\mathcal{S}$ に対して
\begin{equation}
    (PQ)^{-1}S(PQ)=Q^{-1}(P^{-1}SP)Q\in\mathcal{S}
\end{equation}
であるから、積 $PQ$ も再び $N(\mathcal{S})$ に属する。
さらに、$P\in N(\mathcal{S})$ ならば $P^{-1}SP\in\mathcal{S}$ であることから
\begin{equation}
    P(P^{-1}SP)P^{-1}=S\in\mathcal{S}
\end{equation}
が成り立ち、逆元 $P^{-1}$ についても
\begin{equation}
    P S P^{-1}\in\mathcal{S}
\end{equation}
が従う。
したがって、$N(\mathcal{S})$ は一般化パウリ群 $\mathcal{P}_p^n$ の部分群である。

上で見たように、論理演算子はスタビライザー群と可換でなければならない。
実際、一般化パウリ群の元どうしは交換したとき位相因子 $\omega^k$ を生じるだけであるため、
$P$ が全ての $S\in\mathcal{S}$ と可換ならば
\begin{equation}
    P^{-1}SP=S
\end{equation}
となり、$P$ は正規化群に属する。
したがって、論理演算子は正規化群の元として理解することができる。

しかし、正規化群の元のうちスタビライザー群そのものに属するものは、
符号空間上で恒等的に作用するので論理演算子としては自明である。
そこで、符号空間上で非自明に作用する論理演算子は、
正規化群をスタビライザー群で割った剰余類
\begin{equation}
    N(\mathcal{S})/\mathcal{S}
\end{equation}
によって特徴づけられる。
すなわち、2つの正規化群の元 $P,Q\in N(\mathcal{S})$ が
\begin{equation}
    P=QS,\qquad S\in\mathcal{S}
\end{equation}
の形で異なるだけならば、符号空間上では同じ論理演算子を表しているとみなす。
実際、
\begin{equation}
    P\ket{\psi}=QS\ket{\psi}=Q\ket{\psi}\qquad (\ket{\psi}\in\mathcal{C})
\end{equation}
であるから、両者は符号空間上で区別できない。

例えば、qudit 版 Shor符号では
\begin{equation}
    \overline{\mathcal{X}}
    =\mathcal{Z}_1^\dagger \mathcal{Z}_4^\dagger \mathcal{Z}_7^\dagger
\end{equation}
を論理 $\mathcal{X}$ の1つの表現としてとることができた。
ここでブロック内の $\mathcal{Z}$ 型スタビライザー
\begin{equation}
    \mathcal{Z}_1\mathcal{Z}_2^\dagger,\qquad
    \mathcal{Z}_4\mathcal{Z}_5^\dagger,\qquad
    \mathcal{Z}_7\mathcal{Z}_8^\dagger
\end{equation}
を掛けると、
\begin{equation}\begin{aligned}[b]
    &(\mathcal{Z}_1^\dagger \mathcal{Z}_4^\dagger \mathcal{Z}_7^\dagger)
    (\mathcal{Z}_1\mathcal{Z}_2^\dagger)
    (\mathcal{Z}_4\mathcal{Z}_5^\dagger)
    (\mathcal{Z}_7\mathcal{Z}_8^\dagger)\\
    &= \mathcal{Z}_2^\dagger \mathcal{Z}_5^\dagger \mathcal{Z}_8^\dagger
\end{aligned}\end{equation}
となる。
しかし、これらはスタビライザーを掛けただけの違いであるから、符号空間上では同じ論理演算子を表している。
このように、$N(\mathcal{S})/\mathcal{S}$ ではスタビライザー倍だけ異なる演算子を同一視するのであった。

また、先ほど導入した qudit 版 Shor符号の $\overline{\mathcal{X}},\overline{\mathcal{Z}}$ は、
スタビライザー群と可換であり、しかもスタビライザー群そのものには属さない元であった。
したがって、これらは正規化群 $N(\mathcal{S})$ の元であり、
その剰余類 $N(\mathcal{S})/\mathcal{S}$ の中の非自明な元として
論理 qudit 上の一般化パウリ演算子を与えていると解釈できる。

このように、スタビライザー群は符号空間を定め、正規化群の剰余類 $N(\mathcal{S})/\mathcal{S}$ はその中での非自明な論理演算子を与える。
さらに、その中から一般化パウリの交換関係を満たす論理 $\overline{\mathcal{X}},\overline{\mathcal{Z}}$ を選ぶことで、
符号空間の中で論理状態を区別し、それらの間の遷移を表すことができる。

\section{射影演算子とハミルトニアン形式}
ここまでで、スタビライザー群を用いて符号空間という「条件を満たす状態の集合」を代数的に定義できることを見た。
量子回路の立場では、補助量子系を用いてスタビライザー演算子を測定し、
その結果にもとづいて訂正操作やフィードバックを行うことで、
状態を符号空間へと保つという描像をとる。

これに対して物理学の立場では、
この符号空間をある物理的な多体系の基底状態空間として実現するという別の見方が可能になる。
すなわち、外部から逐一測定と訂正を行うアルゴリズム的な描像から離れ、
同じ符号空間を基底状態空間として持つハミルトニアンを構成したいのである。
このような見方をとると、量子誤り訂正符号は「測定と復号のアルゴリズム」としてだけでなく、
多体系の基底状態と励起を持つ物理模型としても理解できる。
特に、スタビライザー条件の破れはハミルトニアンの励起状態として記述されるため、
その励起のエネルギーや局所性を調べることができる。

また、後の章で見るように、このようなハミルトニアンの励起状態を調べることは、エニオンなどの豊かな物理的性質を理解することに直結する。
では、スタビライザー条件 $S_i \ket{\psi} = \ket{\psi}$ を満たす状態のエネルギーを下げるハミルトニアンを、
どのように構成すればよいだろうか。
特に $Z_N$ 系においては、ハミルトニアンの各項の構成方法として
「射影演算子を用いる形式」と「エルミート共役（$\mathrm{h.c.}$）を用いる形式」の2つが自然に現れる。
これらは同じ符号空間を基底状態としながらも、
励起スペクトルに異なる物理的性質をもたらす。

スタビライザー符号では、全てのスタビライザー条件を満たす状態が符号空間をなし、
それらをハミルトニアンの基底状態とみなしたい。
一方、いずれかのスタビライザー条件を破る状態は励起状態に対応してほしい。
したがって、ハミルトニアンの各項は、ある状態が各スタビライザーの $+1$ 固有空間に属しているかどうかを判定するような演算子であるのが自然である。
まず、1つのスタビライザー演算子に対して、その $+1$ 固有空間への射影演算子を構成することを考えよう。
スタビライザー演算子 $S$ に対して、各べき乗を足し合わせた
\begin{equation}
    \Pi_S := \frac{1}{p}\sum_{m=0}^{p-1} S^m
\end{equation}
を定義する。
これは離散フーリエ変換によって固有値 $\omega^k$ のうち $k=0$ の成分だけを取り出す操作に対応している。
実際、$S$ の固有状態 $\ket{\psi}$ が
\begin{equation}
    S\ket{\psi} = \omega^k \ket{\psi},\qquad \omega=e^{2\pi i/p}
\end{equation}
を満たしているとする。
このとき
\begin{equation}\begin{aligned}[b]
    \Pi_S\ket{\psi}
    &= \frac{1}{p}\sum_{m=0}^{p-1}S^m\ket{\psi}
    = \frac{1}{p}\sum_{m=0}^{p-1}\omega^{km}\ket{\psi}
\end{aligned}\end{equation}
となる。
特に $k=0$、すなわち $S\ket{\psi}=\ket{\psi}$ のときには
\begin{equation}
    \Pi_S\ket{\psi}=\frac{1}{p}\sum_{m=0}^{p-1}\ket{\psi}=\ket{\psi}
\end{equation}
である。
一方、$k\neq 0$ のときには、
\begin{equation}
    \Pi_S\ket{\psi}=\frac{1}{p}\sum_{m=0}^{p-1}\omega^{km}\ket{\psi}=\frac{1}{p}\frac{1-\omega^{kp}}{1-\omega^k}\ket{\psi}=0
\end{equation}
となる。
したがって、一般の状態 $\ket{\psi}$ を $S$ の固有空間分解
\begin{equation}
    \ket{\psi}=\sum_{k=0}^{p-1}\ket{\psi^{(k)}},\qquad
    S\ket{\psi^{(k)}}=\omega^k\ket{\psi^{(k)}}
\end{equation}
によって表すと、
\begin{equation}
    \Pi_S\ket{\psi}=\ket{\psi^{(0)}}
\end{equation}
となる。
つまり、$\Pi_S$ は $S$ の $+1$ 固有空間の成分だけを取り出し、
それ以外の成分を消す演算子である。
したがって、$\Pi_S$ は $S$ の $+1$ 固有空間への射影演算子となっている。

この射影演算子を用いると、符号空間を基底状態空間に持つハミルトニアンを定義できる。
最も直接的なものは
\begin{equation}
    H_{\mathrm{proj}}:=-\sum_{i=1}^r \Pi_{S_i}
\end{equation}
である。
ここで $S_1,\dots,S_r$ はスタビライザー群の独立な生成元である。
各項 $\Pi_{S_i}$ は固有値 $0$ または $1$ を持つので、
このハミルトニアンのエネルギーが最小になるのは全ての $i$ に対して
\begin{equation}
    \Pi_{S_i}\ket{\psi}=\ket{\psi}
\end{equation}
が成り立つとき、すなわち
\begin{equation}
    S_i\ket{\psi}=\ket{\psi}
\end{equation}
が全て満たされるときである。
したがって、$H_{\mathrm{proj}}$ の基底状態空間は符号空間 $\mathcal{C}$ に一致する。

ここで注目すべきことに、$H_{\mathrm{proj}}$ は一般には多体相互作用を含むにもかかわらず、
その全ての項が互いに可換である。
例えば、Shor符号においてはスタビライザー演算子は
\begin{equation}\begin{aligned}[b]
    S_{s1}&=\mathcal{Z}_1\mathcal{Z}_2^\dagger&
    S_{p1}&=\mathcal{X}_1\mathcal{X}_2\mathcal{X}_3 \mathcal{X}_4^\dagger\mathcal{X}_5^\dagger\mathcal{X}_6^\dagger
\end{aligned}\end{equation}
であり、
これによる射影演算子は
\begin{equation}\begin{aligned}[b]
    \Pi_{s1} &= \frac{1}{p}\sum_{m=0}^{p-1}(\mathcal{Z}_1\mathcal{Z}_2^\dagger)^m, &
    \Pi_{p1} &= \frac{1}{p}\sum_{m=0}^{p-1}(\mathcal{X}_1\mathcal{X}_2\mathcal{X}_3 \mathcal{X}_4^\dagger\mathcal{X}_5^\dagger\mathcal{X}_6^\dagger)^m
\end{aligned}\end{equation}
といった形をしている。
これを見てわかるように、$\Pi_{s1}$ は2 qudit にのみ作用し、
$\Pi_{p1}$ は6 qudit にのみ作用する項である。
通常、このような多体相互作用を含むハミルトニアンの解析は容易ではない。
しかし、ここでのハミルトニアンはスタビライザー条件から基底状態が明示的にわかっており、
各局所項を同時に対角化できる特別に扱いやすいクラスに属している。
この点は、後に扱うトポロジカル秩序や励起の構造を理解するうえでも重要となる。


\end{document}
